\section{실험 및 결과}
\label{sec:experiments}

본 장에서는 전체 서명 코드의 SW 오류 에뮬레이션으로 상류 연산,
T1, T2 및 RB의 취약 효과를 확인한다. 물리 실험은 T1을 주 대상으로
수행하며, 미수정 융합 구현의 T2 탐색은 소프트웨어 오류모델과 실제
글리치 효과를 비교하기 위한 보조 실험으로 제시한다.

\subsection{실험 환경 및 방법}

실험 장비는 ChipWhisperer-Husky와 CW308 보드에 장착한
STM32F405를 사용하였다. 대상 장치는 HSE 직결 7.37\,MHz로
동작하며, 펌웨어는 \code{HAETAE}-120 전체 서명 과정을 포함한다.
펌웨어 빌드에는 \code{arm-none-eabi-gcc} 13.2.1을 사용하였다.
실험 간 정상 응답과 결함 응답을 비교하기 위해 고정 seed 기반의
결정론 난수발생기를 사용하였다. 이는 오류 효과와 복원식을 검증하기
위한 통제 조건이며 실제 배포용 난수 설정이 아니다.

SW 오류 에뮬레이션에서는 지정한 서명 지점에서 1회만 발화하는 오류
훅을 사용하였다. 상류 지점에서는 seed, 부호, 공개행렬 언패킹 및 LSB
유도값을 변조하였다. T1 훅은 곱 결과를 0으로 대체하고, T2 훅은
응답 계산에서 $\yvec$ 항을 제거한다. RB 훅은 거부 판정을 생략한다.
각 지점은 서로 다른 메시지로 10회씩 세 배치, 총 30회 평가하였다.

물리 실험에서는 글리치 위치(\code{ext\_offset}), 폭, offset 및
repeat를 변화시키며 결과를 수집하였다. 결과는 정상 응답과 동일한
Golden, 장치가 응답하지 않는 Mute, 목표 오류와 다른 값을 생성하는
Other 및 분석 대상인 T1/T2 효과로 구분하였다. 물리 오류의 인과
판정과 계수 비교에는 연구용 펌웨어의 내부 응답과 참 $\svec_1$
스트림을 사용하였다.

\subsection{SW 오류 에뮬레이션 결과}
\label{subsec:software-faults}

Table~\ref{tab:sw-fault-result}은 전체 서명 경로에서 각 의미론적 오류
효과가 방출된 횟수를 나타낸다. 이 비율은 SW 훅을 활성화했을 때 거부
경로를 통과한 빈도이며, 물리 글리치의 성공확률을 의미하지 않는다.

\begin{table}[H]
\centering
\caption{Results of software fault emulation in the full signing path}
\label{tab:sw-fault-result}
\footnotesize
\begin{tabularx}{\columnwidth}{@{}l r X@{}}
\toprule
Point & Observed & Confirmed effect \\
\midrule
SEED    & 30/30 & seed 변조 응답 \\
SIGNBIT & 7/30  & 부호 변조 응답 \\
UNPACK  & 30/30 & 공개행렬 변조 응답 \\
LSB     & 7/30  & challenge 경로 변조 \\
T1      & 7/30  & 곱 항 제거 응답 \\
T2      & 25/30 & nonce 항 제거 및 $\svec_1$ 복원 \\
RB      & 27/30 & 거부 판정 우회 \\
\bottomrule
\end{tabularx}
\end{table}

T2 의미론적 오류가 방출된 25회에는 식~\eqref{eq:t2-recovery}를
적용하였다. 챌린지 $c$와 오류 응답을 NTT 영역으로 변환한 뒤 공개
챌린지의 역원을 곱하여 세 개의 비밀 다항식을 복원하였다. 25회 모두
$\svec_1$의 $3\times256=768$개 계수가 연구용 펌웨어가 출력한 참
값과 일치하였다. 따라서 T2의 관측률은 25/30(83.3\%)이고, 관측된
오류 응답에 대한 계수 복원률은 768/768이었다.

T1 SW 훅에서도 동일 nonce의 정상 응답과 곱 항이 제거된 응답의
차분으로 768개 계수의 복원 관계를 확인하였다. RB는 거부되어야 할
응답이 후속 경로로 진행되는 것을 확인하였으나, 통계적 키 복구는
수행하지 않았다. 따라서 상류 지점과 T1/T2의 결과는 목표 오류 효과와
대수적 영향에 대한 E2 증거이며, RB는 경계 우회의 E2 증거로 한정한다.

\subsection{T1 물리 글리치 및 복원 결과}
\label{subsec:physical-t1}

T1의 물리적 실현 조건을 평가하기 위해 곱 결과 버퍼를 사전에 0으로
초기화한 \code{T1\_CS\_ZEROINIT} 인과대조 구현을 구성하였다.
\code{AXISA\_JIG}는 동일한 nonce와 내부 상태에서 응답 계산을
재실행하여 글리치 파라미터를 반복 평가한다. 정상 실행에서는 세 비밀
다항식 블록이 $c\svec_1$으로 갱신되지만, 글리치로 생략된 반복의
블록은 초기값 0을 유지한다. 따라서 곱셈 반복 생략과 T1 항 제거의
인과관계를 블록 단위로 관찰할 수 있다.

Table~\ref{tab:t1-campaign-result}은 T1 물리 실험의 탐색, 집중 복원
및 별도 재현 과정을 나타낸다. 초기 600회 스윕에서 한 블록의
256/768 복원 구간을 발견하였고, 500회 밀집 탐색에서 동일한
33.3\% 블록 복원이 28회 재현되었다. 전체 응답 구간 1,200회
탐색에서는 약 41.9k, 92.3k 및 143.1k cycle에 세 블록의 고유
글리치 구간이 존재함을 확인하였다.

\begin{table}[H]
\centering
\caption{T1 physical glitch campaigns on the causal-control build}
\label{tab:t1-campaign-result}
\footnotesize
\begin{tabularx}{\columnwidth}{@{}X r X@{}}
\toprule
Campaign & Shots & Main result \\
\midrule
Initial sweep & 600 & 한 블록 구간 식별 \\
Dense reproduction & 500 & 256/768 효과 28회 \\
Full-window scan & 1,200 & 세 블록 구간 식별 \\
Three-band scan & 900 & 최대 256/247/224 \\
Focused recovery & 108/900 & 768/768 복원 \\
Separate run & 594/900 & 768/768 재현 \\
\bottomrule
\end{tabularx}
\end{table}

세 구간을 이용한 첫 900회 탐색에서는 블록별 최대 일치도가
256/256, 247/256 및 224/256으로 측정되었다. 이후 구간을 조정한
집중 실행은 최대 예산 900회 중 108회에서 조기 종료되었다. 이 실행의
\code{ext=41421} 결함 응답은 첫 번째 블록 256개를 복원하였고,
\code{ext=91670} 결함 응답은 나머지 두 블록 512개를 동시에
복원하였다. 두 결함 결과를 누적하여 $\svec_1$ 768/768 계수를
확보하였다.

동일 장비와 빌드에서 수행한 별도 실행은 최대 예산 900회 중 594회에서
동일한 결과를 재현하였다. 이때 \code{ext=41421}에서 첫 블록
256/256, \code{ext=91672}에서 나머지 두 블록 512/512가
복원되었다. 결과적으로 공격 구간을 알고 있는 악용 단계의 최소 출력
집합은 동일 nonce의 정상 기준 응답 1개와 유용한 결함 응답 2개이다.
그러나 구간 탐색, 보정, 첫 복원 및 별도 재현을 모두 포함한 본 연구의
물리 측정량은 총 3,902회이므로, 최소 출력 수를 전체 공격 획득 비용과
동일하게 해석해서는 안 된다.

이 결과는 실제 클럭 글리치가 반복문의 특정 블록을 제거하여 T1의
대수적 관계를 물리적으로 생성할 수 있음을 보인다. 다만 버퍼를
사전 초기화한 인과대조 구현, 고정 nonce 재생 및 내부 응답 판정 조건을
사용하였으므로, 공식 미수정 구현의 종단 공격이 아니라 T1의 구현
의존적 물리 실현성을 보이는 E3 결과로 해석한다.

\subsection{미수정 융합 구현의 물리 T2 비교}
\label{subsec:physical-t2}

미수정 응답 구현은 각 계수에서 난수와 곱 결과를 읽고 더한 뒤 저장하는
융합 구조를 사용한다. 이 구현에 실제 클럭 글리치를 주입하여
$\yvec$ 항만 제거된 깨끗한 T2 응답이 발생하는지 탐색하였다.
1차 스윕 300회에서는 Golden 274회, Mute 21회, Other 5회가
관측되었다. 오류 발생 구간을 중심으로 수행한 350회 특성화에서는
Golden 320회, Mute 25회 및 Other 5회가 관측되었다. 두 실험 모두
깨끗한 T2는 0회였다.

\begin{table}[H]
\centering
\caption{Physical T2 search results on the fused implementation}
\label{tab:physical-t2-result}
\footnotesize
\begin{tabular}{@{}lrrrrr@{}}
\toprule
Test & Trials & Gold. & Mute & Other & T2 \\
\midrule
Sweep & 300 & 274 & 21 & 5 & 0 \\
Characterization & 350 & 320 & 25 & 5 & 0 \\
\bottomrule
\end{tabular}
\end{table}

SW 훅은 벡터 전체의 nonce 항을 한 번에 제거하지만, 융합 구현은
1024개 계수의 load--add--store를 반복한다. 한 번의 국소 글리치로
모든 nonce load만 제거하고 곱 결과와 store를 유지하기는 어렵다.
다만 0/650 결과는 T2의 물리적 불가능성을 증명하지 않으며, 조사한
장치, 빌드, 글리치 방식과 파라미터 범위에서 목표 효과가 관측되지
않았음을 의미한다.

\subsection{결과 비교}

Table~\ref{tab:evidence-result}은 본 연구가 각 취약 지점에 대해
확보한 검증 범위를 요약한다. 본 논문의 핵심 물리 결과는 T1이며,
나머지 지점의 SW 결과는 의미론적 오류가 발생했을 때의 취약 조건과
보안 영향을 확인하는 데 사용한다.

\begin{table*}[t]
\centering
\caption{Scope and evidence of the evaluated HAETAE fault points}
\label{tab:evidence-result}
\small
\begin{tabularx}{0.94\textwidth}{@{}l X X X@{}}
\toprule
Target & SW evidence & Physical evidence & Claim scope \\
\midrule
Upstream &
SEED, SIGNBIT, UNPACK 및 LSB 오류 응답 재현 &
본 연구 범위 밖 &
선행연구 지점의 의미론적 취약성 확인 \\
T2 &
25/30 방출, 각 응답에서 $\svec_1$ 768/768 복원 &
미수정 융합 구현에서 0/650 &
치명적 항 제거 관계와 조사 범위 내 물리 미관측 \\
T1 &
동일 nonce 차분으로 $\svec_1$ 복원 &
인과대조 구현의 두 실행에서 768/768 &
본 연구의 핵심 물리 실현성 결과 \\
RB &
27/30 거부 판정 우회 &
본 연구 범위 밖 &
직접 키 복구가 아닌 잠재적 분포 취약성 \\
\bottomrule
\end{tabularx}
\end{table*}
